\subsection{PROPER MAPPINGS}

If \(f: \mathbf{X} \to \mathbf{Y}\) and \(f': \mathbf{X}' \to \mathbf{Y}'\) are two continuous closed mappings, the product \(f \times f': \mathbf{X} \times \mathbf{X}' \to \mathbf{Y} \times \mathbf{Y}'\) is not necessarily a closed map, even if \(f\) is of the form \(\iota_{\mathbf{X}}\) .

Example. Every constant mapping into a Hausdorff space is closed. But if \(f\) is the constant mapping \(\mathbf{Q} \to 0\) , then \(f \times \iota_{\mathbf{Q}}\) is the mapping \((x, y) \to (0, y)\) of \(\mathbf{Q}^2\) into \(\mathbf{Q}^2\) , so it is the second projection and is not closed (§ 4, no. 2, Remark 1).

DEFINITION 1. Let \(f\) be a mapping of a topological space \(\mathbf{X}\) into a topological space \(\mathbf{Y}\) . \(f\) is said to be proper if \(f\) is continuous and if the mapping \(f \times \iota_{\mathbf{Z}} : \mathbf{X} \times \mathbf{Z} \to \mathbf{Y} \times \mathbf{Z}\) is closed, for every topological space \(\mathbf{Z}\) .

We shall give other characterizations of proper mappings in no. 2 and 3.

If in Definition 1 we take the space Z to consist of a single point, we see that:

PROPOSITION 1. Every proper mapping is closed.

PROPOSITION 2. Let \(f: \mathbf{X} \to \mathbf{Y}\) be a continuous injection. Then the following three statements are equivalent:

\begin{enumerate}
\def\labelenumi{\alph{enumi})}
\item
  \(f\) is proper.
\item
  \(f\) is closed.
\item
  \(f\) is a homeomorphism of \(\mathbf{X}\) onto a closed subset of \(\mathbf{Y}\) .
\end{enumerate}

We have just seen that a) implies b). Since the equivalence relation \(f(x) = f(x')\) is the equality relation, the quotient space of X with respect to this relation can be identified with X; hence b) implies c) by reason of § 5, no. 2, Proposition 3. Finally, if c) is satisfied then \(f \times \iota_{\mathbf{Z}}\) is a homeomorphism of \(\mathbf{X} \times \mathbf{Z}\) onto a closed subspace of \(\mathbf{Y} \times \mathbf{Z}\) and is therefore a closed mapping; hence c) implies a).

PROPOSITION 3. Let \(f: \mathbf{X} \to \mathbf{Y}\) be a continuous mapping. If \(\mathbf{T}\) is any subset of \(\mathbf{Y}\) , let \(f_{\mathbf{T}}\) denote the mapping \(\overline{f}^{1}(\mathbf{T}) \to \mathbf{T}\) which agrees with \(f\) on \(\overline{f}^{1}(\mathbf{T})\) .

\begin{enumerate}
\def\labelenumi{\alph{enumi})}
\item
  If \(f\) is proper, so is \(f_{\mathbf{T}}\) .
\item
  Let \((\mathrm{T}(\iota))_{\iota\in I}\) be a family of subsets of Y whose interiors cover Y, or which is a locally finite closed covering of Y; then if each of the mappings \(f_{\mathbf{T}(\iota)}\) is proper, so is \(f\) .
\end{enumerate}

Let \(\mathbf{Z}\) be a topological space. If \(\mathbf{T}\) is any subset of \(\mathbf{Y}\) , we have

\[
f _ {\mathbf {T}} \times \iota_ {\mathbf {Z}} = (f \times \iota_ {\mathbf {Z}}) _ {\mathbf {T} \times \mathbf {Z}};
\]

if \(f\) is proper, then \(f \times \iota_{\mathbf{Z}}\) is closed, hence so is \((f \times \iota_{\mathbf{Z}})_{\mathbf{T} \times \mathbf{Z}}\) {[}§ 5, no. 1, Proposition 2 a){]}, whence a) is proved. If now \((\mathrm{T}(\iota))_{\iota \in \mathbf{I}}\) satisfies one of the two conditions stated in b), then the covering \((\mathrm{T}(\iota) \times \mathbf{Z})_{\iota \in \mathbf{I}}\) of \(\mathbf{Y} \times \mathbf{Z}\) has the same property; if the \(f_{\mathbf{T}(\iota)}\) are proper then the mappings

\[
(f \times \iota_ {\mathbf {Z}}) _ {\mathbf {T} (\iota) \times \mathbf {Z}}
\]

are closed, whence \(f \times \iota_{\mathbf{Z}}\) is closed {[}§ 5, no. 1, Proposition 2 b){]}. This completes the proof.

PROPOSITION 4. Let I be a finite set and for each \(i \in I\) let \(f_i: X_i \to Y_i\) be a continuous mapping. Let \(X = \prod_{i \in I} X_i\) , \(Y = \prod_{i \in I} Y_i\) , and let \(f: X \to Y\)

be the product mapping \((x_{i})\to (f_{i}(x_{i}))\) . Then:

\begin{enumerate}
\def\labelenumi{\alph{enumi})}
\item
  If each of the \(f_i\) is proper, then \(f\) is proper.
\item
  If \(f\) is proper and if the \(X_i\) are non-empty, then each of the \(f_i\) is proper.
\end{enumerate}

(We shall see in no. 2, Theorem 1, Corollary 3, that this proposition extends to infinite products.)

By induction it is enough to consider the case where \(\mathbf{I} = \{\mathbf{1},\mathbf{2}\}\) . a) Suppose that \(f_{1}, f_{2}\) are proper, and let \(\mathbf{Z}\) be a topological space; \(f_{1} \times f_{2} \times \iota_{\mathbf{Z}}\) is the composition of \(\iota_{\mathbf{Y}_1} \times f_2 \times \iota_{\mathbf{Z}}\) and

\[
f _ {1} \times \iota_ {\mathbf {X} _ {2}} \times \iota_ {\mathbf {Z}};
\]

these two mappings are closed by hypothesis, hence so is \(f_{1} \times f_{2} \times \iota_{\mathbf{Z}}\) {[}§ 5, no. 1, Proposition 1 a){]}, whence \(f_{1} \times f_{2}\) is proper.

\begin{enumerate}
\def\labelenumi{\alph{enumi})}
\setcounter{enumi}{1}
\tightlist
\item
  Now suppose \(f\) is proper. Let \(\mathbf{F}\) be a closed subset of \(\mathbf{X}_2 \times \mathbf{Z}\) and let \(\mathbf{G}\) be the image of \(\mathbf{F}\) in \(\mathbf{Y}_2 \times \mathbf{Z}\) under the mapping \(f_2 \times \iota_{\mathbf{Z}}\) . Then the image of \(\mathbf{X}_1 \times \mathbf{F}\) in \(\mathbf{Y}_1 \times \mathbf{Y}_2 \times \mathbf{Z}\) under \(f_1 \times f_2 \times \iota_{\mathbf{Z}}\) is \(f_1(\mathbf{X}_1) \times \mathbf{G}\) . By hypothesis, this is closed in \(\mathbf{Y}_1 \times \mathbf{Y}_2 \times \mathbf{Z}\) ; if \(\mathbf{X}_1 \neq \emptyset\) , then \(f_1(\mathbf{X}_1)\) is not empty, which implies that \(\mathbf{G}\) is closed in \(\mathbf{Y}_2 \times \mathbf{Z}\) (§ 4, no. 3, Corollary to Proposition 7); hence \(f_2\) is proper. Similarly \(f_1\) is proper if \(\mathbf{X}_2 \neq \emptyset\) .
\end{enumerate}

PROPOSITION 5. Let \(f: \mathbf{X} \to \mathbf{X}'\) and \(g: \mathbf{X}' \to \mathbf{X}''\) be two continuous mappings.

\begin{enumerate}
\def\labelenumi{\alph{enumi})}
\item
  If \(f\) and \(g\) are proper, then \(g \circ f\) is proper.
\item
  If \(g \circ f\) is proper and \(f\) is surjective, then \(g\) is proper.
\item
  If \(g \circ f\) is proper and \(g\) is injective, then \(f\) is proper.
\item
  If \(g \circ f\) is proper and \(X'\) is Hausdorff, then \(f\) is proper.
\end{enumerate}

Let \(\mathbf{Z}\) be a topological space. We have

\[
(g \circ f) \times \iota_ {\mathbf {Z}} = (g \times \iota_ {\mathbf {Z}}) \circ (f \times \iota_ {\mathbf {Z}});
\]

if \(f\) and \(g\) are proper, then \(f \times \iota_{\mathbf{Z}}\) and \(g \times \iota_{\mathbf{Z}}\) are closed; hence {[}§ 5, no. 1, Proposition 1 a){]} ( \(g \circ f\) ) × \(\iota_{\mathbf{Z}}\) is closed; this proves a). The proof of b) {[}resp. c){]} runs along the same lines, using part b) {[}resp. c){]} of Proposition 1 of § 5, no. 1, and remarking that if \(f\) is surjective (resp. if \(g\) is injective) then \(f \times \iota_{\mathbf{Z}}\) is surjective (resp. \(g \times \iota_{\mathbf{Z}}\) is injective). Finally, to prove d), consider the commutative diagram

\[
\mathbf {X} \stackrel {\varphi} {\rightarrow} \mathbf {X} \times \mathbf {X} ^ {\prime}\tag{1}
\]

\[
\begin{array}{c} f \downarrow \\ \mathbf {X} ^ {\prime} \xrightarrow [ \psi ]{} \mathbf {X} ^ {\prime \prime} \end{array} \begin{array}{c} \downarrow (g \circ f) \times \iota_ {\mathbf {x} ^ {\prime}} \\ \times \mathbf {X} ^ {\prime} \end{array}
\]

where \(\varphi(x) = (x, f(x))\) and \(\psi(x') = (g(x'), x')\) . The mapping \(\varphi\) (resp. \(\psi\) ) is a homeomorphism of X (resp. X') onto the graph of \(f\) (resp. the reflection of the graph of \(g\) ) ( \(\S 4\) , no. 1, Proposition 1, Corollary 2). Further, since X' is Hausdorff, the graph \(\varphi(X)\) of \(f\) is closed in X × X' ( \(\S 8\) , no. 1, Proposition 2, Corollary 2). Hence (Proposition 2) \(\varphi\) is proper; on the other hand Proposition 4 shows that ( \(g \circ f\) ) × \(\iota_{X'}\) is proper. By a) above and the commutativity of the diagram (1), \(\psi \circ f\) is proper; but \(\psi\) is injective and therefore \(f\) is proper by c) above.

Remark. If \(\mathbf{X}'\) is not Hausdorff it can happen that \(g \circ f\) is proper and \(f\) not; for example, take \(\mathbf{X}\) and \(\mathbf{X}''\) to consist of one point and \(\mathbf{X}'\) to consist of two points, with the coarsest topology.

COROLLARY I. If \(f: \mathbf{X} \to \mathbf{Y}\) is a proper mapping, then the restriction of \(f\) to a closed subset \(\mathbf{F}\) of \(\mathbf{X}\) is a proper mapping of \(\mathbf{F}\) into \(\mathbf{Y}\) .

For this restriction is the composition \(f \circ j\) , where \(j: \mathbf{F} \to \mathbf{X}\) is the canonical injection, which is proper by Proposition 2.

COROLLARY 2. Let \(f: \mathbf{X} \to \mathbf{Y}\) be a proper mapping, where \(\mathbf{X}\) is Hausdorff. Then the subspace \(f(\mathbf{X})\) of \(\mathbf{Y}\) is Hausdorff.

By reason of Proposition 5 c) we need only consider the case where \(f(\mathbf{X}) = \mathbf{Y}\) . Then the diagonal of \(\mathbf{Y} \times \mathbf{Y}\) is the image under \(f \times f\) of the diagonal of \(\mathbf{X}\) , which is closed (§ 8, no. 1, Proposition 1); \(f \times f\) is proper (Proposition 4); hence the diagonal of \(\mathbf{Y} \times \mathbf{Y}\) is closed (Proposition 1) and therefore \(\mathbf{Y}\) is Hausdorff (§ 8, no. 1, Proposition 1).

COROLLARY 3. Let I be a finite set and for each \(i \in I\) , let \(f_i: X \to Y_i\) be a proper mapping. If X is Hausdorff, then the mapping \(x \to (f_i(x))\) of X into \(\prod_{i \in I} Y_i\) is proper.

This mapping is the composition of the product mapping \((x_{i})\to (f_{i}(x_{i}))\) of \(\mathbf{X}^{\mathbf{I}}\) into \(\prod_{i}\mathbf{Y}_{i}\) and the diagonal mapping of \(\mathbf{X}\) into \(\mathbf{X}^{\mathbf{I}}\) ; since the latter is proper (by Proposition 2 and §8, no. 1, Proposition 1) the conclusion follows from Proposition 4 and Proposition 5 a).

COROLLARY 4. Let X and Y be two topological spaces, \(f: \mathbf{X} \to \mathbf{Y}\) a continuous mapping, R the equivalence relation \(f(x) = f(y)\) on X, and

\[
\mathbf {X} \xrightarrow {p} \mathbf {X} / \mathbf {R} \xrightarrow {h} f (\mathbf {X}) \xrightarrow {i} \mathbf {Y}
\]

the canonical decomposition of \(f\) . Then for \(f\) to be proper it is necessary and sufficient that \(p\) is proper, \(h\) a homeomorphism and \(f(\mathbf{X})\) a closed subset of \(\mathbf{Y}\) .

The conditions are sufficient by virtue of Proposition 5 a) and Proposition 2. Conversely, if \(f\) is proper, then \(f\) is closed; hence \(f(\mathbf{X})\) is closed in \(\mathbf{Y}\) and \(h\) is a homeomorphism (§ 5, no. 2, Proposition 3); also \(h \circ p\) is proper by Proposition 5 c); hence \(p = \overline{h} \circ (h \circ p)\) is proper by Proposition 5 a).

